\section{Introduction}

\subsection{The Question}

A market maker earns the bid--ask spread, but the two sides of a trade do not arrive together.
A dealer who is filled on the bid first is long the stock until someone lifts the ask, and every
minute spent holding that position is exposed to the price. The inventory models of
Garman~\cite{garman1976}, Amihud and Mendelson~\cite{amihudmendelson1980}, Stoll~\cite{stoll1978} and
Ho and Stoll~\cite{hostoll1981} all answer the same question: how should a dealer set prices when
inventory is risky? Avellaneda and Stoikov~\cite{avellaneda2008} restate the answer for a modern
limit order book. Their dealer shifts both quotes away from the side it is already heavy on, and
the size of the shift grows with inventory, risk aversion, volatility and the time left in the
session.

The working paper this project builds on~\cite{amargos2026} derives the model in full and confirms
in a Monte Carlo study that this skew cuts terminal-inventory variance by about an order of
magnitude. Its own scope section is explicit about what it does not do: it \emph{``does not
calibrate or backtest the model against live order-book data''}. That gap is the starting point
here. We ask one question:

\begin{center}
\emph{Does the Avellaneda--Stoikov model survive contact with the National Stock Exchange of India?}
\end{center}

\noindent ``Survive'' has three concrete meanings, and the paper answers them in order.
\begin{enumerate}
\item \textbf{Do its assumptions hold on NSE data?} We test each one and sort it into
  \emph{holds}, \emph{bends} or \emph{breaks}.
\item \textbf{Does its risk control carry over?} We compare the optimal quotes with a symmetric
  dealer quoting the same average spread, first in the model's own world, then on real minutes.
\item \textbf{Does its profit carry over, and is its spread close to the market's?} We split the
  dealer's real PnL into spread capture, adverse selection and inventory carry, and we compare the
  model's spread with the Roll~\cite{roll1984} spread the market actually charges.
\end{enumerate}

\subsection{What We Find}

Figure~\ref{fig:roadmap} summarises the argument. The model's core assumption, that fills
arrive at a rate $Ae^{-k\delta}$ decaying exponentially in the quote distance $\delta$, holds
remarkably well for both stocks. Its picture of the mid-price bends for the illiquid stock,
which shows the bid--ask bounce Roll described, and breaks on intraday volatility, which is
U-shaped rather than constant. In the model's world the optimal quotes beat the symmetric
dealer on every risk measure. On real NSE minutes the risk result survives almost intact, but
the profit result does not: fills on a real exchange carry information, so the price tends to
keep moving against the dealer after each fill, and this adverse selection eats all of the
spread income. The market's own spread is far narrower than the model's, which is what one
expects when a single dealer's optimum is set against a book full of competing liquidity
providers.

\begin{figure}[!htbp]
\centering
\begin{tikzpicture}[
  box/.style={draw=ink2, rounded corners=3pt, align=center, font=\small, minimum height=1.25cm,
              text width=2.75cm, fill=white, line width=0.5pt},
  arr/.style={-{Stealth[length=2.2mm]}, ink2, line width=0.6pt},
  node distance=0.55cm]
  \node[box] (data) {\textbf{NSE data}\\ 1-min bars,\\ liquid + illiquid};
  \node[box, right=of data] (test) {\textbf{Assumptions}\\
    {\color{holds}holds} / {\color{symorange}bends} /\\ {\color{breaks}breaks}};
  \node[box, right=of test] (cal) {\textbf{Calibrate}\\ $\sigma$, $A$, $k$\\ (first \CalibDays{} days)};
  \node[box, right=of cal, yshift=0.85cm] (mc) {\textbf{Model world}\\ Monte Carlo,\\ \Paths{} paths};
  \node[box, right=of cal, yshift=-0.85cm] (rp) {\textbf{Real world}\\ replay on the\\ last 9 days};
  \node[box, below=1.15cm of test, xshift=1.55cm, text width=9.6cm, minimum height=1.0cm, draw=asblue, line width=0.8pt] (verdict)
    {\textbf{Verdict:} inventory control survives; profit does not\\[1pt]
     {\footnotesize adverse selection eats the spread, and the market's spread is narrower than the model's}};
  \draw[arr] (data) -- (test);
  \draw[arr] (test) -- (cal);
  \draw[arr] (cal.east) -- ++(0.25,0) |- (mc.west);
  \draw[arr] (cal.east) -- ++(0.25,0) |- (rp.west);
  \draw[arr] (mc.east) -- ++(0.3,0) |- (verdict.east);
  \draw[arr] (rp.east) -- ++(0.3,0) |- (verdict.east);
\end{tikzpicture}
\caption{The argument of the paper. Each box is a section; each arrow is a result the next step needs.}
\label{fig:roadmap}
\end{figure}

\subsection{Organization}

Section~\ref{sec:model} states the model and the closed-form quotes we test.
Section~\ref{sec:data} describes the data and the two stocks. Section~\ref{sec:assumptions}
tests the assumptions and calibrates the parameters. Section~\ref{sec:nse} shows what the
optimal quotes look like at NSE parameters. Section~\ref{sec:mc} runs the model in its own world
and Section~\ref{sec:replay} runs it on real minutes. Section~\ref{sec:discussion} compares the
two and sets the model's spread against the market's, and Section~\ref{sec:conclusion}
concludes. The appendix documents the code and data needed to reproduce every figure.
